%%! Author = lili
%%! Date = 6/15/26

\chapter{Tutorium 15.06.2026}\label{ch:tutorium-15.06.2026}
\untit{Präsenzaufgaben 2.2: Grenzwertsätze \hfill 15.--19.\ Juni 2026}
\txo{Bearbeiten Sie auf jeden Fall Aufgabe 2.2.4. Falls Ihnen Zeit bleibt, bearbeiten Sie
zusätzlich die weiteren Aufgaben.}
Binomialverteilung benutzt man, wenn es aus der Aufgabenstellungen hervor sinnvoll ist oder wenn die Aufgabenstellung
sagt, dass etwas binomialverteilt ist.
\begin{align*}
    \pe(X_n \leq \delta) & = \pe(X_n - np \leq \delta - np) \\
    & \pe(\frac{X_n - np}{\sqrt{npq}} \leq \frac{\delta - np}{\sqrt{npq}}) \\
    & \overset{*}{\approx} \phi(\frac{\delta - np }{\sqrt{npq}})
\end{align*}
* das ist das Anwenden von Moire-Laplace.

Tscheby:
\[
    \pe(|X- \E{X} | \geq \delta) \leq \frac{Var(X)}{\delta^2}
\]

\section{PÜ 2.2.4. (Lösung da, noch nicht eingetragen)}\label{sec:pu-2.2.4.}
% TODO Bilder Nils
Der Lufthansa Airbus A380 bietet insgesamt 526 Fluggästen Platz.
Da Kunden manchmal
ihren Flug nicht antreten, lassen Fluggesellschaften zwecks optimaler Auslastung Überbuchungen
zu.
Es sollen möglichst viele Tickets verkauft werden, wobei jedoch die Wahrscheinlichkeit
einer Überbuchung maximal 5\% betragen soll.
Erfahrungsgemäß erscheinen Kunden mit einer
Wahrscheinlichkeit von 10\% nicht zum Flug.

\begin{auf}
    \subsection*{(a) Geben Sie einen geeigneten Wahrscheinlichkeitsraum an, der es erlaubt, die folgenden
    Aufgabenteile zu bearbeiten. Diskutieren Sie, welche Annahmen Sie machen müssen, um Ihre Wahl des Wahrschein-
    lichkeitsraums zu rechtfertigen.}
    \index{Aufgabe!Wahrscheinlichkeitsraum angeben}
    \index{Aufgabe!Annahmen}
    \paragraph{Lösung Tutor}
    \[
        \eraum = \{ \underbrace{0}_{kommt\ nicht}, \underbrace{1}_{kommt}\}^n
    \]
    kommt nicht ist $q = \frac{1}{10}$ und kommt ist $p = \frac{9}{10}$.
    \[
        \pe( \{\ele\}) = p^{\sum_{i=1}^n \ele_i} \cdot (1-p)^{\sum_{i=1}^n \ele_i}
    \]
    Bspw. $\{\ele\} = (1,0,1,1,0)$

    Wenn Binomialverteilt, geht man davon aus jeder hat die gleich WS zu kommen und sie sind unabhängig (wir
    ignorieren bspw., dass Familien meistens zusammen reisen oder absagen).
\end{auf}

%%Der Erwartungswert ist nicht die WS aller möglichen Kombinationen, sondern die WS gewichtet nach de...
Erwartungswert berechnet den Durchschnitt einer ZV!

\begin{auf}
    \subsubsection*{(b) Geben Sie ein Verfahren an, wie die Frage, wie viele Tickets unter den oben genannten
    Bedingungen maximal angeboten werden können, exakt beantwortet werden könnte}
    \paragraph{Lösung des Tutors:} Man definiert eine ZV $S_n$ über die aufsummiert wird:
    \[
        S_n( \underbrace{\ele}_{\ele = \omega_1, ... \omega_n}) = \sum^n_{i=1} \ele_i
    \]
    $S_n$ ist die Anzahl der Leute, die kommen. \\
    Wenn wir die WS der Überbuchung auf maximal 0,05 halten wollen, interessiert uns
    \[
        \pe( S_n > 526 ) \leq 0,05
    \]
    \begin{itemize}
        \item $n$ = Anzahl verkaufte Tickets (wir wollen n maximieren)
        \item $S_n$ = Anzahl der Personen, die kommen
    \end{itemize}
    Auf Basis dessen oben:
    \[
        \underbrace{\pe(S_n \leq 526) }_{ = \sum^{526}_{i=1} \binom{n}{i} p^i (1-p)^{n-i}} \geq 0,95
    \]
    Starte bei n = 526 und berechne $\pe(S_n \leq 526)$ mit inkrementierendem n und wenn es bei einem n über 0,95
    liegt, dann ist das das gesuchte n das n-1 ist.
    \begin{align*}
        \pe(S_n \leq 526) & = \pe(S_n - np \leq 526 - np) \\
        & = \pe(\frac{S_n - np}{\sqrt{npq}} \leq \frac{256 - np}{\sqrt{npq}}) \\
        & \approx \phi(\underbrace{\frac{526 - np }{\sqrt{npq}}}_{\geq 1,65})
    \end{align*}
    \[
        \frac{526 - np }{\sqrt{npq}} \geq 1,65
    \]
    p und q einsetzen
    \[
        526 - (\sqrt{n})^2 \cdot \frac{9}{10} \geq 1,65 \sqrt{n} \cdot \frac{3}{10}
    \]
    und dann landet man durch Rumschieben bei
    \[
        ( \underbrace{\sqrt{n}}_{x} )^2  \underbrace{+ 1,65 \cdot \frac{1}{3}}_{c_1} \underbrace{\sqrt{n}}_{x}  \underbrace{- \frac{10}{q o. 9?}}_{c_2} \leq 0
    \]
    % TODO
    TODO noch das letzte von der Tafel
\end{auf}

\begin{auf}
    \subsubsection*{(c) Nutzen Sie den Satz von de Moivre-Laplace, um näherungsweise zu bestimmen, wie viele
    Tickets unter den oben genannten Bedingungen maximal angeboten werden können.}
    \index{Aufgabe!Moivre-Laplace}
\end{auf}


\section{PÜ 2.2.5. (Lösung da, noch nicht eingetragen)}\label{sec:pu-2.2.5}
% TODO Bilder Nils
Eine Süßwarenfirma gibt an, daß 25\% ihrer Gummibärchen die beliebte Farbe rot haben.
Im
folgenden gehen wir davon aus, daß die Angaben des Herstellers richtig sind.
Wir nehmen auch
an, daß die Gummibärchen unabhängig voneinander rot sind oder eine andere Farbe haben.

Trotzdem waren in einer kleineren Stichprobe lediglich 20\% der Gummibärchen rot.
Wir wollen
der Sache auf den Grund gehen und kaufen daher n weitere Gummibärchen.
Sei $b_n$ die
Wahrscheinlichkeit, dass höchstens 20\% der neu gekauften n Gummibärchen rot sind.

\begin{auf}
    \subsection*{(a) Geben Sie für jedes beliebige $n \in \nz$ einen geeigneten Wahrscheinlichkeitsraum an, der es
    erlaubt, die folgenden Aufgabenteile zu bearbeiten}
    \index{Aufgabe!Wahrscheinlichkeitsraum angeben}
    \[
        \eraum = \{
        \underbrace{0}_{nicht\ rot}, \underbrace{1}_{rot}
        \}^n
    \]
    wobei $n$ die Anzahl der gekauften Gummibärchen ist.
    \[
        \pe( \{\ele\}) = p^{\sum_{i=1}^n \ele_i} \cdot (1-p)^{\sum_{i=1}^n \ele_i}
    \]
    mit p=0,25 = WS für rot und q = 1-p
\end{auf}

\begin{auf}
    \subsubsection*{(b) Definieren Sie eine Zufallsvariable $S_n$, die die Anzahl der roten Gummibärchen unter den
    n gekauften Gummibärchen angibt.}
    \index{Aufgabe!Zufallsvariable definieren}
    \[
        S_n = \sum_{i=1}^n \ele_i, \quad S_n \sim Bin(n,p)
    \]
\end{auf}

\begin{auf}
    \subsubsection*{(c) Sei n = 20. Berechnen Sie $b_{20}$ auf drei Nachkommastellen genau, d.h., geben Sie die
    Wahrscheinlichkeit in der Form 0.xyz bzw. xy.z \% an.}
    \index{Aufgabe!Verteilung berechnen!Binomialverteilung}
    \[
        \frac{S_n}{n} (=mitHütchenzeichen) Durchschnittliche\ Anzahl\ rot
    \]
    Es gilt
    \[
        \pe( \frac{1}{n} S_n \leq \frac{1}{5} ) = \pe(S_n \leq \frac{n}{5})
    \]
    Und mit n = 20
    \[
        \pe(S_{20} \leq \frac{20}{5}) = \pe(S_{20} \leq 4) = \sum_{i=1}^n \binom{20}{i} p^i \cdot (1-p)^{20-i} = 0,4148...
    \]
\end{auf}

\begin{auf}
    \subsubsection*{(d) Beweisen Sie mit Hilfe der Tschebyscheffschen Ungleichung, daß $b_{200} < b_{20}$ gilt.}
    \index{Aufgabe!Tschebyscheffsche Ungleichung}
    Also $b_n = \pe(\frac{1}{n} S_n \leq \frac{1}{5})$ \\
    Heißt wir müssen $ \pe(|X- \E{X} | \geq \delta) \leq \frac{Var(X)}{\delta^2}$ in jene Form umwandeln. Erstmal
    Erwartungswert von beiden Seiten abziehen(?):
    \[
        \pe(\frac{1}{n} S_n - \E{\frac{1}{n} S_n} \leq \frac{1}{5} - \E{\frac{1}{n} S_n})
    \]
    Es ist ja
    \[
        \E{\frac{1}{n} S_n} = - \frac{1}{n} \E{S_n} = \frac{1}{n} \cdot np = p \overset{Aufgabe}{=} \frac{1}{4}
    \]
    Also
    \begin{align*}
        \pe(\frac{1}{n} S_n - \E{\frac{1}{n} S_n} \leq \frac{1}{5} - \E{\frac{1}{n} S_n}) & =
        \pe(\frac{1}{n} S_n - \E{\frac{1}{n} S_n} \leq - \frac{1}{20}) \\
        & =  \pe(| \frac{1}{n} S_n - \E{\frac{1}{n} S_n} | \geq \frac{1}{20})
    \end{align*}
    % TODO
    TODO die Sachen von der Tafel noch zu ende
\end{auf}

\begin{auf}
    \subsubsection*{(e) Zeigen Sie mit Hilfe der Tschebyscheffschen Ungleichung, dass
        $\lim_{n \rightarrow \infty} b_n = 0$ gilt}
    \index{Aufgabe!Tschebyscheffsche Ungleichung}
    NOCH KEINE LÖSUNG % TODO NOCH KEINE LÖSUNG
\end{auf}

\section{PÜ 2.2.6. (Lösung da, noch nicht eingetragen)}
% TODO Bilder Nils
Drei faire Würfel werden unabhängig voneinander geworfen. Die Zufallsvariable $Z$
bezeichne deren Augensumme.

\begin{auf}[Geben Sie einen geeigneten Wahrscheinlichkeitsraum an für dieses
Zufallsexperiment, und definieren Sie $Z$ formal als Abbildung auf dem
zugehörigen Ereignisraum.]
    \index{Aufgabe!Wahrscheinlichkeitsraum angeben}
    NOCH KEINE LÖSUNG % TODO NOCH KEINE LÖSUNG
\end{auf}

\begin{auf}[Berechnen Sie die Wahrscheinlichkeit, dass die Augensumme der drei Würfel
gleich $5$ ist.]
    \index{Aufgabe!Wahrscheinlichkeit berechnen}
    NOCH KEINE LÖSUNG % TODO NOCH KEINE LÖSUNG
\end{auf}

Wir wiederholen nun dieses Zufallsexperiment, bis die Augensumme zum ersten Mal gleich
$5$ ist. Die Wiederholungen seien unabhängig. Die Zufallsvariable $\tau$ beschreibe die
Nummer des Wurfs, bei dem die Augensumme zum ersten Mal gleich $5$ ist.


\begin{auf}[Geben Sie einen geeigneten Wahrscheinlichkeitsraum an für dieses
Zufallsexperiment, und definieren Sie $\tau$ formal als Abbildung auf dem
zugehörigen Ereignisraum.]
    \index{Aufgabe!Wahrscheinlichkeitsraum angeben}
    NOCH KEINE LÖSUNG % TODO NOCH KEINE LÖSUNG
\end{auf}

\begin{auf}[Welche Verteilung hat die Zufallsvariable $\tau$? Geben Sie, falls
erforderlich, die Parameter der Verteilung von $\tau$ an.]
    \index{Aufgabe!Verteilung!Typ bestimmen}
    NOCH KEINE LÖSUNG % TODO NOCH KEINE LÖSUNG
\end{auf}

\begin{auf}[Bestimmen Sie den Erwartungswert und die Varianz von $\tau$.]
    NOCH KEINE LÖSUNG % TODO NOCH KEINE LÖSUNG
    \index{Aufgabe!Varianz!Varianz berechnen}\index{Aufgabe!Erwartungswert!Erwartungswert berechnen}
\end{auf}

\begin{auf}[(f) Sei $k \in \nz \setminus \{1\}$. Bestimmen Sie mit Hilfe der Tschebyscheffschen
Ungleichung eine obere Schranke für die Wahrscheinlichkeit, dass $\tau$ einen
Wert annimmt, der mindestens dem $k$-fachen Erwartungswert von $\tau$
entspricht.]
    \index{Aufgabe!Tschebyscheffsche Ungleichung}\index{Aufgabe!Wahrscheinlichkeit berechnen}
    NOCH KEINE LÖSUNG % TODO NOCH KEINE LÖSUNG
\end{auf}

\begin{auf}[Bestimmen Sie für $k = 2$ die Wahrscheinlichkeit aus (f) mit Hilfe des
Computers exakt, und vergleichen Sie mit dem Ergebnis der Abschätzung in (f).]
    NOCH KEINE LÖSUNG % TODO NOCH KEINE LÖSUNG
\end{auf}


%\section{Besprechung ÜB 2.1.1}\label{sec:besprechung-ub-2.1.1}
%Am Rande des Leinewebermarkts wird Ihnen folgendes Spiel vorgeschlagen.
%Es werden zwei faire
%Würfel geworfen.
%Anschließend wird das Maximum und das Minimum der beiden Ergebnisse
%gebildet.
%Der Einsatz des Spiels beträgt zwei Euro, und Sie erhalten die Di!erenz von Maximum
%und Minimum ausgezahlt.

%\begin{auf}
%    \subsection*{(a) Geben Sie einen Wahrscheinlichkeitsraum \wraum an, der es erlaubt, die folgenden Auf-
%    gabenteile zu bearbeiten}
%\end{auf}

%\begin{auf}
%    \subsubsection*{(b) Definieren Sie auf dem von Ihnen gewählten Wahrscheinlichkeitsraum \wraum eine Zu-
%    fallsvariable M, die das Maximum der Ergebnisse der beiden Würfe angibt, und eine
%    Zufallsvariable N, die das Minimum der Ergebnisse der beiden Würfe angibt)}
%\end{auf}

%\begin{auf}
%    \subsubsection*{(c) Die Zufallsvariable
%        $G : \eraum \rightarrow \rz$ bezeichne Ihren Gewinn, sollten Sie eine Runde spielen.
%    Dabei sei G negativ, falls Sie Geld verlieren. Definieren Sie diese Zufallsvariable mit Hilfe
%    der Zufallsvariablen M und N , und geben Sie den Bildbereich $G(\eraum)$ an.}
%\end{auf}

%\begin{auf}
%    \subsubsection*{(d) Berechnen Sie den erwarteten Gewinn, d.h. den Erwartungswert von G}
%\end{auf}

%\begin{auf}
%    \subsubsection*{(e) Ist das Spiel fair , das heißt, beträgt der erwartete Gewinn 0 Euro?}
%\end{auf}

%\section{Besprechung ÜB 2.1.2}\label{sec:besprechung-ub-2.1.2}
%Die gemeinsame Dichte der Zufallsvariablen X und Y sei gegeben durch
%\[
%    f (x, y) = \begin{cases}
%                   ce^{-2(x+y)},& f\ddot ur\ 0 \leq x \leq y \geq \infty, \\
%                   0, & sonst
%    \end{cases}
%\]

%\begin{auf}
%    \subsection*{(a) Bestimmen Sie c derart, daß f eine \gls{zweidimensionaledichte} ist}
%\end{auf}

%\begin{auf}
%    \subsubsection*{(b) Berechnen Sie Erwartungswert und Varianz von X und Y . Warum existieren diese?}
%\end{auf}

%\begin{auf}
%    \subsubsection*{(c) Berechnen Sie den Erwartungswert von 2X - Y . Warum existiert dieser?}
%\end{auf}

%\section{Besprechung ÜB 2.1.3}\label{sec:besprechung-ub-2.1.3}
%\begin{auf}
%    \subsection*{(a)}
%\end{auf}
%\begin{auf}
%    \subsubsection*{(b)}
%\end{auf}
%\begin{auf}
%    \subsubsection*{(c)}
%\end{auf}
%\begin{auf}
%    \subsubsection*{(d)}
%\end{auf}
%\begin{auf}
%    \subsubsection*{(e)}
%\end{auf}
%\begin{auf}
%    \subsubsection*{(f)}
%\end{auf}
%\begin{auf}
%    \subsubsection*{(g)}
%\end{auf}